% GENERATED FILE. Edit canonical lessons, then run npm run generate:books.
% canonical-source-hash: fnv1a64:854c793d5e1838c4
% canonical-lessons: TE-C394-vivarana, TE-C394-spastata, TE-C394-samarthana, TE-C394-samarthincu, TE-C394-khaccitamaina

\chapter{Explaining and Justifying}
\label{ch:te-explaining-and-justifying}
\hlchaptermodality{394}

\begin{chapteropening}
\textbf{By the end of this chapter:} \emph{I can explain, clarify and justify a view.}

Complete the last lesson of chapter 394: I can explain, clarify and justify a view.
\end{chapteropening}

\section[\texorpdfstring{\te{వివరణ} (vivaraṇa)}{వివరణ (vivaraṇa)}]{\te{వివరణ} (vivaraṇa) — an explanation; a clarification}

\label{lesson:TE-C394-vivarana}



\begin{quote}
Before the new one: say the Telugu for harmful, then the Telugu for to be useful.
\end{quote}

\subsection*{You'll want to know: \te{వివరణ}}
\textbf{\te{వివరణ}} — \emph{vivaraṇa} — \textquotedblleft{}an explanation; a clarification\textquotedblright{}.

\begin{quote}
\te{అతని} \te{వివరణ} \te{నాకు} \te{అర్థం} \te{కాలేదు}.
\end{quote}

\begin{quote}
\emph{atani vivaraṇa nāku arthaṁ kālēdu.}
\end{quote}

\begin{quote}
I did not understand his explanation.
\end{quote}

The noun of \te{వివరించు}, to explain, which you know.

\begin{cousinweb}
\textbf{\te{వివరణ}} (\emph{vivaraṇa}) is borrowed from Sanskrit \emph{vivaraṇa}, \textquotedblleft{}an uncovering, an opening\textquotedblright{}, and so \textquotedblleft{}explanation, exposition\textquotedblright{} (Monier-Williams).
\end{cousinweb}

\begin{grammarlens}[title={What you've built}]
One more word for giving reasons and opinions.
\end{grammarlens}

\subsection*{Your turn}
\begin{itemize}
\raggedright
  \item \emph{Say it:} \emph{vivaraṇa}
  \item \emph{Say it:} \emph{vivaraṇa}, once more
  \item \emph{Say it:} \emph{atani vivaraṇa nāku arthaṁ kālēdu.} — I did not understand his explanation.
  \item \emph{Recall:} say the Telugu for harmful, then the Telugu for to be useful, then say \emph{vivaraṇa} again
\end{itemize}

\begin{tcolorbox}[breakable,colback=teal!4,colframe=teal!35!black,title={Before you move on}]
What does \te{వివరణ} mean? (\textquotedblleft{}An explanation; a clarification\textquotedblright{}.) Say it once more.
\end{tcolorbox}
\section[\texorpdfstring{\te{స్పష్టత} (spaṣṭata)}{స్పష్టత (spaṣṭata)}]{\te{స్పష్టత} (spaṣṭata) — clarity}

\label{lesson:TE-C394-spastata}



\begin{quote}
Before the new one: say the Telugu for to be useful, then the Telugu for an explanation.
\end{quote}

\subsection*{You'll want to know: \te{స్పష్టత}}
\textbf{\te{స్పష్టత}} — \emph{spaṣṭata} — \textquotedblleft{}clarity\textquotedblright{}.

\begin{quote}
\te{ఆ} \te{నియమం} \te{గురించి} \te{ఇంకా} \te{స్పష్టత} \te{రాలేదు}.
\end{quote}

\begin{quote}
\emph{ā niyamaṁ guriñci iṅkā spaṣṭata rālēdu.}
\end{quote}

\begin{quote}
There is still no clarity about that rule.
\end{quote}

\te{స్పష్టం}, clear, + -\te{త}: the same family as \te{స్పష్టమైన}, which you know.

\begin{cousinweb}
\textbf{\te{స్పష్టత}} (\emph{spaṣṭata}) is built on Sanskrit \emph{spaṣṭa}, \textquotedblleft{}clearly perceived, distinct, plain\textquotedblright{} (Monier-Williams), with the abstract ending -\te{త}.
\end{cousinweb}

\begin{grammarlens}[title={What you've built}]
One more word for giving reasons and opinions.
\end{grammarlens}

\subsection*{Your turn}
\begin{itemize}
\raggedright
  \item \emph{Say it:} \emph{spaṣṭata}
  \item \emph{Say it:} \emph{spaṣṭata}, once more
  \item \emph{Say it:} \emph{ā niyamaṁ guriñci iṅkā spaṣṭata rālēdu.} — There is still no clarity about that rule.
  \item \emph{Recall:} say the Telugu for to be useful, then the Telugu for an explanation, then say \emph{spaṣṭata} again
\end{itemize}

\begin{tcolorbox}[breakable,colback=teal!4,colframe=teal!35!black,title={Before you move on}]
What does \te{స్పష్టత} mean? (\textquotedblleft{}Clarity\textquotedblright{}.) Say it once more.
\end{tcolorbox}
\section[\texorpdfstring{\te{సమర్థన} (samarthana)}{సమర్థన (samarthana)}]{\te{సమర్థన} (samarthana) — a justification; a defence}

\label{lesson:TE-C394-samarthana}



\begin{quote}
Before the new one: say the Telugu for an explanation, then the Telugu for clarity.
\end{quote}

\subsection*{You'll want to know: \te{సమర్థన}}
\textbf{\te{సమర్థన}} — \emph{samarthana} — \textquotedblleft{}a justification; a defence\textquotedblright{}.

\begin{quote}
\te{తన} \te{తప్పుకు} \te{అతని} \te{దగ్గర} \te{ఏ} \te{సమర్థన} \te{లేదు}.
\end{quote}

\begin{quote}
\emph{tana tappuku atani daggara ē samarthana lēdu.}
\end{quote}

\begin{quote}
He has no justification for his mistake.
\end{quote}

The reasons that back an act or a view.

\begin{cousinweb}
\textbf{\te{సమర్థన}} (\emph{samarthana}) is borrowed from Sanskrit \emph{samarthana}, \textquotedblleft{}reflection, deliberation\textquotedblright{} (Monier-Williams); Telugu uses it for the case made for an act or a view.
\end{cousinweb}

\begin{grammarlens}[title={What you've built}]
One more word for giving reasons and opinions.
\end{grammarlens}

\subsection*{Your turn}
\begin{itemize}
\raggedright
  \item \emph{Say it:} \emph{samarthana}
  \item \emph{Say it:} \emph{samarthana}, once more
  \item \emph{Say it:} \emph{tana tappuku atani daggara ē samarthana lēdu.} — He has no justification for his mistake.
  \item \emph{Recall:} say the Telugu for an explanation, then the Telugu for clarity, then say \emph{samarthana} again
\end{itemize}

\begin{tcolorbox}[breakable,colback=teal!4,colframe=teal!35!black,title={Before you move on}]
What does \te{సమర్థన} mean? (\textquotedblleft{}A justification; a defence\textquotedblright{}.) Say it once more.
\end{tcolorbox}
\section[\texorpdfstring{\te{సమర్థించు} (samarthiñcu)}{సమర్థించు (samarthiñcu)}]{\te{సమర్థించు} (samarthiñcu) — to justify; to defend (a view)}

\label{lesson:TE-C394-samarthincu}



\begin{quote}
Before the new one: say the Telugu for clarity, then the Telugu for a justification.
\end{quote}

\subsection*{You'll want to know: \te{సమర్థించు}}
\textbf{\te{సమర్థించు}} — \emph{samarthiñcu} — \textquotedblleft{}to justify; to defend (a view)\textquotedblright{}.

\begin{quote}
\te{అతను} \te{తన} \te{పనిని} \te{గట్టిగా} \te{సమర్థించాడు}.
\end{quote}

\begin{quote}
\emph{atanu tana panini gaṭṭigā samarthiñcāḍu.}
\end{quote}

\begin{quote}
He firmly defended his work.
\end{quote}

Past (I …): \te{సమర్థించాను}, \emph{samarthiñcānu}.

Also to back someone else: \te{ఆ} \te{ప్రతిపాదనను} \te{సమర్థించు}, to support that proposal.

\begin{cousinweb}
\textbf{\te{సమర్థించు}} (\emph{samarthiñcu}) is built on Sanskrit \emph{samarthana} (Monier-Williams), with the Telugu verb ending -\te{ఇంచు}; the noun is \te{సమర్థన}.
\end{cousinweb}

\begin{grammarlens}[title={What you've built}]
One more word for giving reasons and opinions.
\end{grammarlens}

\subsection*{Your turn}
\begin{itemize}
\raggedright
  \item \emph{Say it:} \emph{samarthiñcu}
  \item \emph{Say it:} \emph{samarthiñcu}, once more
  \item \emph{Say it:} \emph{atanu tana panini gaṭṭigā samarthiñcāḍu.} — He firmly defended his work.
  \item \emph{Recall:} say the Telugu for clarity, then the Telugu for a justification, then say \emph{samarthiñcu} again
\end{itemize}

\begin{tcolorbox}[breakable,colback=teal!4,colframe=teal!35!black,title={Before you move on}]
What does \te{సమర్థించు} mean? (\textquotedblleft{}To justify; to defend (a view)\textquotedblright{}.) Say it once more.
\end{tcolorbox}
\section[\texorpdfstring{\te{ఖచ్చితమైన} (khaccitamaina)}{ఖచ్చితమైన (khaccitamaina)}]{\te{ఖచ్చితమైన} (khaccitamaina) — exact; precise}

\label{lesson:TE-C394-khaccitamaina}



\begin{quote}
Before the new one: say the Telugu for a justification, then the Telugu for to justify.
\end{quote}

\subsection*{You'll want to know: \te{ఖచ్చితమైన}}
\textbf{\te{ఖచ్చితమైన}} — \emph{khaccitamaina} — \textquotedblleft{}exact; precise\textquotedblright{}.

\begin{quote}
\te{అతను} \te{ఖచ్చితమైన} \te{లెక్క} \te{చెప్పాడు}.
\end{quote}

\begin{quote}
\emph{atanu khaccitamaina lekka ceppāḍu.}
\end{quote}

\begin{quote}
He gave the exact figure.
\end{quote}

\te{ఖచ్చితంగా}, exactly, for certain.

\begin{cousinweb}
\textbf{\te{ఖచ్చితమైన}} (\emph{khaccitamaina}) has an origin that is not settled: it may go back to Sanskrit \emph{khacita}, \textquotedblleft{}set, inlaid\textquotedblright{} (Monier-Williams), but the step to \textquotedblleft{}exact\textquotedblright{} is unclear. The ending -\te{మైన} makes an adjective.
\end{cousinweb}

\begin{grammarlens}[title={What you've built}]
One more word for giving reasons and opinions.
\end{grammarlens}

\subsection*{Your turn}
\begin{itemize}
\raggedright
  \item \emph{Say it:} \emph{khaccitamaina}
  \item \emph{Say it:} \emph{khaccitamaina}, once more
  \item \emph{Say it:} \emph{atanu khaccitamaina lekka ceppāḍu.} — He gave the exact figure.
  \item \emph{Recall:} say the Telugu for a justification, then the Telugu for to justify, then say \emph{khaccitamaina} again
\end{itemize}

\begin{tcolorbox}[breakable,colback=teal!4,colframe=teal!35!black,title={Before you move on}]
What does \te{ఖచ్చితమైన} mean? (\textquotedblleft{}Exact; precise\textquotedblright{}.) Say it once more.
\end{tcolorbox}
